\section{Over-the-Air Computation: Foundations and Techniques}
\subsection{Theoretical Foundations of OAC}
\begin{itemize}
    \item \textit{Signal Superposition in the Multiple Access Channel (MAC)}:
        The distinct feature of OAC where computation is handled by harnessing
        interference via simultaneous transmissions
        \cite{sahin_survey_2023,goldenbaum_robust_2013}.
    \item \textit{Nomographic Functions}: Defining the class of
        functions computable via
        OAC, which can be expressed as a post-processed sum of pre-processed
        sensor readings \cite{sahin_survey_2023,goldenbaum_analyzing_2011,
        goldenbaum_nomographic_2015,goldenbaum_robust_2013}.
    \item \textit{Examples of Nomographic Functions}: Arithmetic
        mean, weighted sum,
        maximum/minimum value, and modulo-2
        sum\cite{sahin_survey_2023,evgenidis_waveform_2025,goldenbaum_robust_2013}.
\end{itemize}
\subsection{OAC Schemes and Implementation Strategies}

    \subsubsection{Nomographic Functions}
    The nomographic functions are the functions that depict the behavior of the \gls{ota} schemes.
    In this thesis project I will focus on how different schemes can be used to compute the arithmetic mean of the transmitted data.
    

\begin{itemize}
    \item \textit{Analog vs. Digital OAC}: Reviewing uncoded analog
        transmission (often
        optimal for simple Gaussian sensor networks) versus digital schemes
        designed for modern systems (e.g., using QAM or one-bit
        quantization)\cite{evgenidis_waveform_2025,zhu_one-bit_2021}.
    \item \textit{Coherent Schemes (CSI Dependent)}: Discussing
        methods requiring
        Channel State Information (CSI), such as Truncated Channel
        Inversion (TCI) or
        Zero-Forcing (ZF) coordination, for phase and amplitude
        alignment\cite{sahin_survey_2023,cao_optimized_2022,zhu_one-bit_2021}.
    \item \textit{Non-Coherent Schemes (CSI Independent)}: Examining
        techniques robust to
        synchronization errors, such as those using Orthogonal Signaling
        (FSK/PPM)\cite{sahin_survey_2023,yang_federated_2020}.
\end{itemize}

    In this section I will demonstrate the Raised Cosine method in order to take advantage of
    the signal superposition in the \gls{mac}. More specifically I will implement three devices
    transmitting simultaneously to a fusion center. The transmitted data will be pre-processed 
    using the Raised Cosine method in order to mitigate the intersymbol interference (ISI) that 
    is inherent in digital communication transceivers. The flowgraph of the implementation is 
    shown in Fig. \ref{fig:gnuradio_flowgraph}. In it three devices pick a character from a
    sequence of charachters, unpack it to separate bytes and perform the mapping to the 
    constellation points. The resulting symbols are then passed through a Raised Cosine filter 
    and transmitted simultaneously to the fusion center. The fusion center performs the matched 
    filtering and the necessary post-processing to compute the desired function of the transmitted 
    data.
    \begin{landscape}
        \begin{figure}[htpb]
            \centering
            \includegraphics[width=\linewidth]{ota_isi_matched.pdf}
            \caption{GNU Radio Companion flowgraph for the OTA implementation including ISI considerations.}
            \label{fig:gnuradio_flowgraph}
        \end{figure}
    \end{landscape}
\subsection{Applications Context (e.g., Distributed Consensus)}
    Application choice (e.g., Federated Edge Learning (FEEL) for gradient
    aggregation, or average consensus in (Wireless Sensor Networks (WSNs))
    
    \cite{sahin_survey_2023,yang_distributed_2025,yang_federated_2020,
    cao_optimized_2022,zhu_over--air_2024}.
